% English translation, made from our own transcription (original.tex), not from any existing
% modern translation. All mathematics, equation numbers (\tag), \origpage markers, and figures
% are preserved unchanged; only the prose is translated. "Schediasma" is rendered "Sketch"
% (matching the fagnano-1718-lemniscata translation); "arc." (the author's abbreviation for
% arco/arc) is kept. The three probable misprints noted in original.tex's header ("Teorema IV"
% printed twice, for what should be "VI"; "equazioni (6), e (8)" where "(7), e (8)" is meant; and
% "dz" for a likely "dx" in Scolio I's "Terzo" remark) are reproduced as printed, unresolved, per
% the same transcription.
\documentclass{article}
\usepackage{readmasters}
\begin{document}

\origpage{304}
\section*{XXXIV. Method for measuring the lemniscate}
\subsection*{Sketch II (*)}
\emph{(*) Giornale de' letterati d'Italia, vol.~XXX, p.~87.}

\textbf{Theorem I.} --- Let the two equations (1) and (2) written below be given; I say that,
the first of them being posited, the other also holds

\[ x = \frac{\sqrt{1 \mp \sqrt{1-z^{4}}}}{z}. \tag{1} \]

\[ \frac{\pm dz}{\sqrt{1-z^{4}}} = \frac{dx\sqrt{2}}{\sqrt{1+x^{4}}}. \tag{2} \]

\textbf{Proof.} --- Equation (1), differentiated, furnishes a value of $dx$ which, being reduced
to a simple expression, gives

\[ dx = \frac{\pm dz\sqrt{1 \mp \sqrt{1-z^{4}}}}{z^{2}\sqrt{1-z^{4}}}. \]

From the same equation (1) one further deduces the following

\[ \frac{\sqrt{1+x^{4}}}{\sqrt{2}} = \frac{\sqrt{1 \mp \sqrt{1-z^{4}}}}{z^{2}}. \]

And dividing the second-to-last equation by this last one, one arrives at equation (2).

Which was to be demonstrated.

\textbf{Theorem II.} --- Let the two equations (3) and (4) written below be given; I say that,
the first of it being posited, the other also holds

\[ x = \frac{\sqrt{1\mp z}}{\sqrt{1\pm z}}. \tag{3} \]

\[ \frac{\mp dz}{\sqrt{1-z^{4}}} = \frac{dx\sqrt{2}}{\sqrt{1+x^{4}}}. \tag{4} \]

\origpage{305}
\textbf{Proof.} --- Differentiating equation (3), and proceeding in the proper manner, one has
$dx = \dfrac{\pm dz}{\sqrt{1-z^{2}}} \times \dfrac{1}{1\pm z}$.

Moreover the same equation (3) shows $\dfrac{\sqrt{1+x^{4}}}{\sqrt{2}} = \dfrac{\sqrt{1+z^{2}}}{1\pm z}$.

And dividing the second-to-last equation by this last one, one obtains equation (4).

Which was to be demonstrated.

\textbf{Theorem III.} --- Let the two equations (5) and (6) written below be given; I say that,
the first of them being posited, the other also holds

\[ x = \frac{u\sqrt{2}}{\sqrt{1-u^{4}}}. \tag{5} \]

\[ \frac{du}{\sqrt{1-u^{4}}} = \frac{1}{\sqrt{2}}\times\frac{dx}{\sqrt{1+x^{4}}}. \tag{6} \]

\textbf{Proof.} --- From equation (5), differentiated and duly handled, there results
$\dfrac{dx}{\sqrt{2}} = \dfrac{du}{\sqrt{1-u^{4}}} \times \dfrac{1+u^{4}}{1-u^{4}}$, and from the
said equation (5) one deduces $\sqrt{1+x^{4}} = \dfrac{1+u^{4}}{1-u^{4}}$.

Dividing thereafter the second-to-last equation by this last one, one arrives at equation (6).

Which was to be demonstrated.

\textbf{Theorem IV.} --- Let the two equations (6) and (8) written below be given; I say that,
the first of them being posited, the other also holds

\[ x = \frac{\sqrt{1-t^{4}}}{t\sqrt{2}}. \tag{7} \]

\[ \frac{-dt}{\sqrt{1-t^{4}}} = \frac{1}{\sqrt{2}}\times\frac{dx}{\sqrt{1+x^{4}}}. \tag{8} \]

\textbf{Proof.} --- Differentiating equation (7), and reducing the value of $dx$ to a convenient
expression, one finds $dx\sqrt{2} = \dfrac{-dt}{\sqrt{1-t^{4}}} \times \dfrac{1+t^{4}}{t^{2}}$.

The same equation (7) shows $2\sqrt{1+x^{4}} = \dfrac{1+t^{4}}{t^{2}}$.

\origpage{306}
And the second-to-last equation, divided by this last one, yields equation (8).

Which was to be demonstrated.

\textbf{Scholium I.} --- To make use of these theorems let it be noted:

First, that $\displaystyle\int \frac{dz}{\sqrt{1-z^{4}}}$ represents the direct arc $CS$ of the
lemniscate (fig.~29) subtended by the chord $CS = z$, provided the semi-axis $CL$ of this curve
is equal to unity.

\rmfigure{figures/fig-29.png}{Fig.~29.}{Fig. 29: the lemniscate with chord $CS=z$ and its
semi-axis $CL$, from the original plate (Tav. III, Opere Matematiche vol. II).}

Second, that $\displaystyle\int \frac{-dz}{\sqrt{1-z^{4}}}$ expresses the inverse arc $SL$ of the
same curve, which corresponds to the aforesaid chord $CS$ ($z$).

Third, that $\displaystyle\int \frac{dz}{\sqrt{1+x^{4}}} = \frac{3}{2}\int dx\,\sqrt{1+x^{4}} -
\frac{1}{2}\sqrt{1+x^{4}}$, as is proved by calculation.

Fourth, that $\displaystyle\int dx\,\sqrt{1+x^{4}}$ expresses the arc $AQ$ (fig.~30) of the
primary cubic parabola corresponding to the abscissa $AF = x$, provided that, taking $a$ as
unity, the parameter of this parabola is $a\sqrt{3}$, and the ordinates are normal to the
abscissa.

\rmfigure{figures/fig-30.png}{Fig.~30.}{Fig. 30: the primary cubic parabola with abscissa
$AF=x$, ordinate $FQ$, tangent portion $QV$, and origin $A$, from the original plate (Tav. III,
Opere Matematiche vol. II).}

Fifth, that $x\sqrt{1+x^{4}}$ denotes the portion $QV$ of the tangent (fig.~30) contained
between the ordinate $FQ$ and the indefinite line $AV$, which from the origin $A$ of the
parabola runs parallel to the ordinates.

Sixth, that in consequence one will have

\[ \int \frac{dx}{\sqrt{1+x^{4}}} = \frac{3}{2}\,\mathrm{arc.}\,AQ - \frac{1}{2}\,QV. \]

\emph{Examples to show the way of applying the preceding theorems to geometry.}

\emph{Example I, for Theorem I} (fig.~29 and 30). --- In equation (1) let the arbitrary sign
$\mp$ signify the negative sign; let equation (2) be integrated by means of the preceding
scholium, and one will find

\[ \mathrm{Arc.}\,CS = \frac{3}{\sqrt{2}}\,\mathrm{arc.}\,AQ - \frac{1}{\sqrt{2}}\,QV. \]

Annulling $AF\,(x)$ will also annul $CS\,(z)$, and so this equation is complete.

\emph{Example II, for Theorem II} (fig.~29 and 30). --- In equation (3) let the sign $\mp$
represent the negative sign; let equation (4) be integrated, and one will disc-
\origpage{307}
over

\[ \mathrm{Arc.}\,SL = \frac{3}{\sqrt{2}}\,\mathrm{arc.}\,AQ - \frac{1}{\sqrt{2}}\,QV. \]

The supposition $CS\,(z) = 1$, which annihilates the inverse arc $SL$, renders $AF\,(x) = 0$,
and so this equation is complete.

\textbf{Scholium II.} --- In the course of this paper I shall no longer prove that the equations
are complete, it being enough for readers to have the method by which I have proceeded in these
two examples, so that they may satisfy themselves of the completeness of the following
equations; they will also be able to make use of the four preceding theorems to draw out other
ways of measuring the lemniscate by means of the extent of the primary cubic parabola, not
omitting, where it is needed, the subtraction or addition of the constant quantity proper to
render the equations complete.

They will moreover be able to deduce from these theorems truths altogether new, concerning the
comparison of the arcs of the said parabola, etc.

I shall content myself with observing that, since the measure of the primary cubic parabola
depends in several ways on the extent of the lemniscate by virtue of the four preceding
theorems, and the measure of the lemniscate depends on the extent of the equilateral hyperbola
and of a kind of ellipse \emph{jointly}, as I discovered in Sketch I, it follows that the measure
of the primary cubic parabola depends in several ways on the extent of the two aforesaid conic
sections \emph{jointly}. This invention cannot fail to please whoever shall have considered what
may be read in the Acta of Leipzig of the year 1695, p.~64, and p.~184 toward the end.

\textbf{Theorem V.} --- Let the two equations (9) and (10) written below be given; I say that,
the first of them being posited, the other also holds

\[ \frac{u\sqrt{2}}{\sqrt{1-u^{4}}} = \frac{1}{z}\sqrt{1-\sqrt{1-z^{4}}}. \tag{9} \]

\[ \frac{dz}{\sqrt{1-z^{4}}} = \frac{2\,du}{\sqrt{1-u^{4}}}. \tag{10} \]

\textbf{Proof.} --- Setting in equation (1) of Theorem I the negative sign in place of $\mp$, and
in place of $x$ its value $\dfrac{u\sqrt{2}}{\sqrt{1-u^{4}}}$ noted in equation (5) of Theorem
III, one has equation (9); then sur-
\origpage{308}
rogating in equation (2) of Theorem I, in place of $\dfrac{dx\sqrt{2}}{\sqrt{1+x^{4}}}$, its
value $\dfrac{2\,du}{\sqrt{1-u^{4}}}$, which arises from the supposition made of $x$ (as is seen
by multiplying equation (6) of Theorem III by 2), equation (10) follows.

Which was to be demonstrated.

\textbf{Corollary I} (fig.~29). --- Let $z$ be called the chord $CS$ of the lemniscate, and $u$
the other chord $CI$, and integrating equation (10), one will find

\[ \mathrm{Arc.}\,CS = 2\,\mathrm{arc.}\,CI. \]

Whence, considering the letter $z$ as known in equation (9), one will draw from it the value of
the chord $CI\,(u)$ that bisects the direct arc $CS$.

\textbf{Corollary II} (fig.~29). --- From equation (9) arises this other

\[ z = \frac{2u\sqrt{1-u^{4}}}{1+u^{4}}. \tag{11} \]

And this value of $z$ is such that, calling $u$ the chord $CI$ taken as known, and making the
other chord $CS$ equal to the second member of equation (11), the direct arc $CS$ is double the
other direct arc $CI$.

\textbf{Theorem IV.} --- Let the two equations (12) and (13) written below be given; I say that,
the first of them being posited, the other also holds

\[ \frac{\sqrt{1-t^{4}}}{t\sqrt{2}} = \frac{1}{z}\sqrt{1-\sqrt{1-z^{4}}}. \tag{12} \]

\[ \frac{dz}{\sqrt{1-z^{4}}} = -\frac{2\,dt}{\sqrt{1-t^{4}}}. \tag{13} \]

\textbf{Proof.} --- Setting in equation (1) of Theorem I the negative sign in place of $\mp$, and
in place of $x$ its value $\dfrac{1}{t\sqrt{2}}\sqrt{1-t^{4}}$ noted in equation (7) of Theorem
IV, one has equation (12). Then substi-
\origpage{309}
tuting in equation (2) of Theorem I, in place of $\dfrac{dx\sqrt{2}}{\sqrt{1+x^{4}}}$, its value
$\dfrac{-2\,dt}{\sqrt{1-t^{4}}}$, which derives from the supposition made of $x$ (as appears by
multiplying equation (8) of Theorem IV by 2), one obtains equation (13).

Which was to be demonstrated.

\textbf{Corollary I} (fig.~29). --- Let $z$ be named the chord $CS$ of the lemniscate, and ($t$)
the other chord $CH$, and there will result arc. $CS = 2$ arc. $HL$.

\textbf{Corollary II} (fig.~29). --- Rearranging equation (12), and putting in it $u$ in place of
$z$, and $r$ in place of $t$, one arrives at this other

\[ u = \frac{2r\sqrt{1-r^{4}}}{1+r^{4}}. \tag{14} \]

So that, naming $r$ the chord $CO$, and taking the other chord $CI\,(u)$ equal to the second
member of equation (14), the direct arc $CI$ will be double the inverse arc $OL$.

\textbf{Problem I.} --- To cut the quadrant of the lemniscate into three equal parts.

\textbf{Solution} (fig.~29 and 31). --- By virtue of Theorem VI and its Corollary I, the direct
arc $CS$ (fig.~29) will be double the inverse arc $HL$, provided that, calling $t$ the chord
$CH$, and $z$ the other chord $CS$, one has equation (12). Suppose now that the two points $S$
and $H$ coincide in $T$; this supposition will render $CS\,(z) = CH\,(t)$, and the direct arc
$CT$ (fig.~31) will be double the inverse arc $TL$, which will become the third part of the
quadrant.

\rmfigure{figures/fig-31.png}{Fig.~31.}{Fig. 31: the lemniscate's quadrant with chord $CI$ and
points $T$, $L$ used in the trisection and quintisection problems, from the original plate
(Tav. III, Opere Matematiche vol. II).}

But supposing $t = z$ in equation (12), and thereafter arranging it in the proper manner, etc.,
one obtains

\[ z = \sqrt[4]{-3+2\sqrt{3}}. \]

Whence, attributing to the chord $CT$ (fig.~31) this value of $z$, and taking the direct arc $CI$
equal to the inverse arc $TL$ by means of Theorem III of Sketch I, one will have the solution of
the problem.

Which was to be found.

\textbf{Problem II.} --- To cut the quadrant of the lemniscate into five equal parts.
\origpage{310}
\textbf{Solution} (fig.~29 and 31). --- Let $u$ be called the chord $CI$ (fig.~29) of the
lemniscate, and $z$ the other chord $CS$; let $z = \dfrac{2u\sqrt{1-u^{4}}}{1+u^{4}}$, and the
arc $CS$ will be double the arc $CI$ by the second corollary of Theorem V. Let $r$ then be named
the chord $CO$ (fig.~29); suppose the chord $CI\,(u) = \dfrac{2r\sqrt{1-r^{4}}}{1+r^{4}}$, and
the direct arc $CI$ will be double the inverse arc $OL$ by Corollary II of Theorem VI; therefore
the direct arc $CS$ will be quadruple the inverse arc $OL$.

Suppose now $CS\,(z) = CO\,(r)$; the points $S$ and $O$ will coincide in $T$, the arc $TL$
(fig.~32) will be one fifth of the quadrant, the arc $CI = 2$ arc. $TL$ will contain two fifths
of the same quadrant, and the arc $IT = $ arc. $CI$ will contain the other two fifths of it.

Let there then be taken, by means of Theorem III of Sketch I, the inverse arc $EL$ (fig.~32)
equal to the direct arc $CI$, and the direct arc $CB$ equal to the inverse arc $TL$, and the
points $B$, $I$, $E$, $T$ will cut the quadrant into five equal parts.

It remains only to find the precise value of the chord $CT$ (fig.~32), equal at the same time to
$z$ and to $r$, and this will be found by substituting in equation (14), in place of $r$, the
letter $z$, whence one will have

\[ u = \frac{2z\sqrt{1-z^{4}}}{1+z^{4}}. \tag{15} \]

And putting this value of $u$ in equation (11), one will arrive at another equation which will
contain no unknown but the letter $z$, and which must be reckoned an equation of the eighth
degree, since if in it one supposes $z^{4} = b$, there will result a new equation of exactly the
eighth degree; the calculation will be longer than difficult, provided one does not neglect
those ways of facilitating it which are well known to skilled analysts, and so I judge it
superfluous to set it out in this paper. I shall only remark that, in the equation expressing the
value of $z^{4}$, the same quantity $z^{4}$ will have two true values less than $a^{4}$; the
greater of these two values will express the fourth power of $CT\,(z)$, and the lesser of them
will express the fourth power of $CI\,(u)$; for if in equation (15) one puts, in place of $z$,
its value in $u$ drawn from equation (11), one will find an equation altogether similar to that
which expresses the value of $z^{4}$, etc.

Which was to be found.

The late Monsieur the Marquis de l'Hospital, in his excellent \emph{Treatise on Conic}
\origpage{311}
\emph{Sections}, book 9, prop.\ 9, teaches a way of constructing equations of the eighth degree.

\textbf{Theorem VII.} --- Let the two equations (16) and (17) written below be given; I say that,
the first of them being posited, the other also holds

\[ \frac{u\sqrt{2}}{\sqrt{1-u^{4}}} = \frac{\sqrt{1-z}}{\sqrt{1+z}}. \tag{16} \]

\[ \frac{-dz}{\sqrt{1-z^{4}}} = \frac{2\,du}{\sqrt{1-u^{4}}}. \tag{17} \]

Which was to be demonstrated.

\textbf{Proof.} --- In equation (3) of Theorem II let the sign $\mp$ represent the negative sign,
and in place of $x$ let there be put its value $\dfrac{u\sqrt{2}}{\sqrt{1-u^{4}}}$ noted in
equation (5) of Theorem III, and one will have equation (16). Let there then be surrogated in
equation (4) of Theorem II, in place of $\dfrac{dx\sqrt{2}}{\sqrt{1+x^{4}}}$, its value
$\dfrac{2\,du}{\sqrt{1-u^{4}}}$, which comes from the supposition made of $x$ (as equation (6) of
Theorem III multiplied by 2 shows), and one will have equation (17).

Which was to be demonstrated.

\textbf{Corollary} (fig.~29). --- Calling $z$ the chord $CS$ of the lemniscate, and $u$ the other
chord $CI$, and thereafter integrating equation (17), one obtains Arc. $SL = 2$ arc. $CI$.

Provided that $u$ and $z$ have the value that is deduced from equation (16).

\textbf{Problem III} (fig.~33). --- Let the quadrant of the lemniscate be bisected at $M$; and
let the direct arc $CB$, less than half the quadrant, be given: to find the intermediate arc $MI$
equal to the given direct arc $CB$.

\rmfigure{figures/fig-33.png}{Fig.~33.}{Fig. 33: the lemniscate's quadrant with points $B$, $I$,
$M$, $E$ used in the bisection-transfer construction, from the original plate (Tav. III,
Opere Matematiche vol. II).}

\textbf{Solution.} --- I. Let there be taken, by means of Theorem VII and its corollary, the
inverse arc $EL$ double the given direct arc $CB$. II. Let there be taken, by means of Theorem V
and its Corollary I, the direct arc $CI$ subduple of the direct arc $CE$. It is evident that the
sum of the arcs $CB + CI$,
\origpage{312}
being equal to half the sum of the arcs $CE + EL$, will consequently be equal to half the
quadrant; therefore

\[ \mathrm{arc.}\,CB + \mathrm{arc.}\,CI = \mathrm{arc.}\,CM, \]

and taking away from both members of this equation the common arc $CI$, there remains arc.
$CB = $ arc. $MI$.

Which was to be demonstrated.

\textbf{Corollary} (fig.~33). --- Let there be taken, by means of Theorem III of Sketch I, the
inverse arc $FL$ equal to the direct arc $CI$, and one will have

\[ \mathrm{arc.}\,Mf = \mathrm{arc.}\,MI = \mathrm{arc.}\,CB. \]

\textbf{Theorem VIII.} --- Let the two equations (18) and (19) written below be given; I say
that, the first of them being posited, the other also holds

\[ \frac{\sqrt{1-t^{4}}}{t\sqrt{2}} = \frac{\sqrt{1-z}}{\sqrt{1+z}}. \tag{18} \]

\[ \frac{-dz}{\sqrt{1-z^{4}}} = -\frac{2\,dt}{\sqrt{1-t^{4}}}. \tag{19} \]

\textbf{Proof.} --- In equation (3) of Theorem II let the sign $\mp$ represent the negative sign,
and in place of $x$ let there be put its value $\dfrac{\sqrt{1-t^{4}}}{t\sqrt{2}}$ expressed in
equation (7) of Theorem IV, and one will have equation (18). Let there then be substituted in
equation (4) of Theorem II, in place of $\dfrac{dx\sqrt{2}}{\sqrt{1+x^{4}}}$, its value
$-\dfrac{2\,dt}{\sqrt{1-t^{4}}}$, which results from the supposition made of $x$ (as equation (8)
of Theorem IV multiplied by 2 shows), and one will arrive at equation (19).

\textbf{Corollary} (fig.~29). --- Naming $z$ the chord $CS$ of the lemniscate, and ($t$) the
other chord $CH$, attributing to $z$ and to $t$ their due values deduced from equation (18), and
integrating equation (19), one discovers the inverse arc $SL$ equal to double the inverse arc
$HL$.

\textbf{Scholium III.} --- From Theorems V, VI, VII, and VIII, and their corollaries, readers
will be able to deduce ways of bisecting the quadrant of the lemniscate entirely uniform with
the manner of doing so discovered by me
\origpage{313}
in Sketch I, and this will serve as proof of the correctness and fruitfulness of my method.

\textbf{Problem IV.} --- Given that the quadrant of the lemniscate is cut into a given number of
equal parts, to subdivide each of the said parts into two equally equal portions.

\textbf{Solution.} --- For greater clarity and brevity let those direct arcs of the lemniscate
that contain an odd number of its equal parts be called odd direct arcs. By virtue of Corollary I
of Theorem V let the direct arcs subduple of all the odd arcs be taken; then, by Corollary I of
Theorem III of Sketch I, let the inverse arcs equal to all these subduple direct arcs be found,
and the intent will be obtained.

Which was to be found.

\textbf{Scholium IV.} --- From Theorem VIII readers will be able to deduce another, similar way
of resolving this problem, which I omit for brevity.

\textbf{Corollary.} --- Since in Sketch I I have taught the way of cutting the quadrant of the
lemniscate into two equal parts, and in Problems I and II of the present paper I have found the
way of cutting the same quadrant into three and five equal parts, it follows, by virtue of the
solution of this problem, that the quadrant of the lemniscate can be divided algebraically into
as many equal parts as are contained in these three formulas, namely: $2 \times 2^{m}$;
$3 \times 2^{m}$; $5 \times 2^{m}$, in which the exponent $m$ signifies any positive whole number.

And this is a new and singular property of my curve.

\end{document}
